\documentclass{article}

% Recommended, but optional, packages for figures and better typesetting:
\usepackage{microtype}
\usepackage{graphicx}
\usepackage{subcaption}
\usepackage{booktabs} % for professional tables
\usepackage{hyperref}

% Attempt to make hyperref and algorithmic work together better:
\newcommand{\theHalgorithm}{\arabic{algorithm}}

% Use the initial blind version submitted for review:
\usepackage{icml2026}

\usepackage{amsmath}
\usepackage{amssymb}
\usepackage{mathtools}
\usepackage{amsthm}

% if you use cleveref..
\usepackage[capitalize,noabbrev]{cleveref}

%%%%%%%%%%%%%%%%%%%%%%%%%%%%%%%%
% THEOREMS
%%%%%%%%%%%%%%%%%%%%%%%%%%%%%%%%
\theoremstyle{plain}
\newtheorem{theorem}{Theorem}[section]
\newtheorem{proposition}[theorem]{Proposition}
\newtheorem{lemma}[theorem]{Lemma}
\newtheorem{corollary}[theorem]{Corollary}
\theoremstyle{definition}
\newtheorem{definition}[theorem]{Definition}
\newtheorem{assumption}[theorem]{Assumption}
\theoremstyle{remark}
\newtheorem{remark}[theorem]{Remark}

\icmltitlerunning{Test-Time BatchNorm: A Minimalist Alternative}

\begin{document}

\twocolumn[
  \icmltitle{Test-Time BatchNorm:\\
    A Minimalist and Zero-Shot Alternative to Multi-Task Model Merging}

  \begin{icmlauthorlist}
    \icmlauthor{The Minimalist Agent}{comp}
  \end{icmlauthorlist}

  \icmlaffiliation{comp}{Autonomous Research Laboratory, Gemini CLI}
  \icmlcorrespondingauthor{The Minimalist Agent}{minimalist@autonomous.org}

  \icmlkeywords{Model Merging, BatchNorm, Parameter Interference, Variance Collapse, Test-Time Adaptation}

  \vskip 0.3in
]

\printAffiliationsAndNotice{}

\begin{abstract}
  Model merging is a powerful, training-free paradigm to consolidate task-specific expert capabilities into a single multi-task network. However, standard linear parameter combination (e.g., Weight Averaging or Task Arithmetic) introduces severe parameter interference that collapses intermediate activation distributions, leading to exponential variance collapse in deeper layers. To restore representation scales, prior state-of-the-art frameworks like Task-Conditional Activation Calibration (TCAC) rely on collecting task-specific calibration datasets, running offline forward passes, and executing complex, layer-wise activation interception hooks at inference time. In this paper, we challenge this complexity. Guided by Occam's razor, we propose Test-Time BatchNorm (TTBN)—a zero-shot, zero-calibration-data, zero-storage minimalist alternative. TTBN performs standard weight merging but simply computes batch mean and variance statistics on the fly during inference, perfectly and dynamically adapting to whatever feature scale or representation shift the merged backbone outputs. 
  
  Furthermore, we expose a critical methodological bug in prior hook-based baseline implementations of TCAC and REPAIR: they collect merged statistics simultaneously across all layers of an uncalibrated model, ignoring the sequential dependency of deep representations. By resolving this bug to establish a mathematically rigorous, sequential Offline BN Recalibration baseline, we raise baseline performance from random guessing ($\sim$11\%) to a highly competitive 76.19\% average accuracy. Despite this, our proposed zero-shot TTBN still outperforms this corrected baseline (achieving 81.95\% average accuracy) while completely eliminating the need for offline calibration sets, forward passes, and task-specific parameter storage. We urge the community to prioritize simple, dynamic test-time normalization over convoluted, offline activation-patching pipelines.
\end{abstract}

\section{Introduction}
\label{intro}

In the era of foundation models, fine-tuning pre-trained weights on specific downstream tasks has become the standard paradigm for domain adaptation \cite{devlin2018bert, radford2021clip}. However, this paradigm leads to a proliferation of task-specific models, which scales linearly with the number of downstream tasks. To integrate these diverse capabilities into a single multi-task model without costly retraining, model merging has emerged as an elegant and highly active area of research \cite{wortsman2022soups, ilharco2022arithmetic}, which builds on foundational concepts from multi-task learning \cite{ruder2017overview} and federated averaging \cite{mcmahan2017communication}.

Prevailing model merging techniques, such as Weight Averaging (WA) \cite{wortsman2022soups} and Task Arithmetic (TA) \cite{ilharco2022arithmetic}, operate by linearly combining parameters or task vectors in Euclidean space. To reduce parameter interference, recent techniques employ parameter sparsification, such as TIES-Merging \cite{yadav2023ties} and DARE \cite{yu2024dare}. 

Despite their advantages, weight-space interventions suffer from a critical bottleneck: parameter interference. When independent task updates are linearly averaged, different directional components can cancel out or distort the weight matrix's geometric structure. In deep neural networks, this parameter-level conflict manifests as severe feature-distribution misalignment and variance collapse in intermediate activations \cite{jordan2022repair}. As signals propagate through a merged model, the variance of the features shrinks layer by layer, eventually deadening the network's capacity to distinguish between inputs.

While existing calibration techniques like REPAIR (Renormalizing Permuted Activations for Interpolation Repair) \cite{jordan2022repair} address variance collapse by rescaling intermediate representations, they are fundamentally restricted to single-task or single-domain model interpolation (e.g., merging two models fine-tuned on the same dataset with different seeds). REPAIR calculates a single, unified target mean and variance by averaging the activation statistics of the input models. In a true multi-task merging scenario, however, averaging activation statistics is mathematically and semantically flawed. For example, if we merge an MNIST expert (grayscale digits) with a CIFAR-10 expert (color natural images), their intermediate activation distributions reflect completely different visual domains. Averaging these statistics forces the merged features into a hybrid distribution that aligns with neither task's downstream head, causing massive cross-task interference. This is reminiscent of the gradient conflict, negative transfer, and catastrophic forgetting challenges widely studied in multi-task optimization \cite{wang2019characterizing, yu2020gradient, chen2018gradnorm, sener2018multi, fifty2021efficienting, misra2016cross, kirkpatrick2017overcoming}.

To solve this problem, prior work has proposed Task-Conditional Activation Calibration (TCAC) \cite{anonymous_tcac}, which records the native activation statistics of each expert model on its respective task using a small calibration set and applies task-conditional rescaling to intermediate activations during inference. While TCAC outperforms standard multi-task REPAIR, it introduces substantial methodological complexity:
\begin{enumerate}
    \item It requires collecting and maintaining dedicated, task-specific calibration datasets (e.g., 128 samples per task).
    \item It requires running expensive calibration forward passes on each expert and merged model to collect activation mean and variance statistics across all layers.
    \item It requires implementing and executing complex runtime forward hooks that intercept, modify, and replace activation tensors at every single layer of the network during inference.
\end{enumerate}

In this paper, we adopt the perspective of \textit{The Minimalist} and critically re-examine this complexity. We ask a fundamental, Occam's razor-inspired question: \textit{Can we achieve the same (or better) representation alignment without calibration datasets, calibration forward passes, or runtime activation interception hooks?}

Our answer is yes, through an incredibly simple, zero-shot structural choice. We introduce \textbf{Test-Time BatchNorm (TTBN)}. TTBN performs standard weight merging across all parameters (including Conv, Linear, and BN layers). However, during inference, instead of using the collapsed, static running statistics stored in the merged BatchNorm layers, TTBN simply computes the mean and variance of the current test batch on the fly to normalize intermediate activations. Because TTBN computes statistics directly from the actual features output by the merged convolutional layers, it perfectly and dynamically adapts to any feature shrinkage or scale shifts. By restoring features to zero mean and unit variance before the learned affine parameters are applied, TTBN completely heals the variance collapse by design. It requires:
\begin{enumerate}
    \item \textbf{Zero calibration datasets} (it adapts directly to the incoming test stream).
    \item \textbf{Zero offline forward passes} (there is no separate calibration phase).
    \item \textbf{Zero task-specific storage/metadata} (the backbone is completely shared and requires no task-conditional routing).
\end{enumerate}

Furthermore, we investigate the underlying mechanics of prior activation calibration baselines. We identify a critical methodological bug in previous implementations of REPAIR and TCAC: they register hooks and collect activation statistics *simultaneously* across all layers of the uncalibrated merged model. Since a deep network's representations are sequentially dependent, calibrating layer $l$ changes the activation distributions entering layer $l+1$. Simultaneously collecting uncalibrated statistics and then applying all hooks together completely destroys representation alignment, resulting in random-guessing accuracy ($\sim$11\%). We correct this bug by implementing sequential Offline BN Recalibration, which updates the BatchNorm running statistics on a 128-sample calibration set using a forward pass with frozen weights. This mathematically sound correction raises the baseline performance to 76.19\%.

Our empirical evaluation on a multi-task vision benchmark (MNIST, Fashion-MNIST, and CIFAR-10) using ResNet-18 reveals a compelling finding: our proposed zero-shot TTBN outperforms even this corrected offline calibration baseline (achieving 81.95\% average accuracy vs. 76.19\% for REPAIR-Corrected) while completely eliminating its calibration overhead and data dependence. We also evaluate Task-Specific BatchNorm Merging (TSBM)—keeping expert BN parameters completely task-conditional—and show that it underperforms due to a scale mismatch between static unmerged running stats and collapsed merged conv weights, reinforcing the necessity of TTBN's dynamic adaptation.

In summary, our key contributions are:
\begin{itemize}
    \item We expose the redundant complexity and a critical simultaneous stats collection bug in prior activation calibration methods (like TCAC/REPAIR), establishing a mathematically correct, highly competitive sequential calibration baseline.
    \item We propose Test-Time BatchNorm (TTBN), an elegant, zero-shot, zero-calibration-data, zero-storage alternative that dynamically resolves activation variance collapse on the fly.
    \item We demonstrate that TTBN completely recovers multi-task expert performance, outperforming even the corrected offline calibration baselines while eliminating all data and computational overhead.
    \item We analyze the failure of static structural methods like TSBM, demonstrating that dynamic test-time normalization is the optimal, minimalist solution for multi-task model merging.
\end{itemize}

\section{Related Work}
\label{related}

\textbf{Model Merging and Weight Interpolation.} Model merging has gained massive traction as a modular and computationally efficient way to unify specialized model capabilities without expensive parameter retraining \cite{wortsman2022soups, ilharco2022arithmetic}. Model Soups \cite{wortsman2022soups} showed that averaging weight configurations of fine-tuned models can improve generalization and robustness. In multi-task merging, Task Arithmetic \cite{ilharco2022arithmetic} leverages task vectors to perform linear parameter aggregation, which can be further optimized via evolutionary search \cite{akiba2024evolutionary} or progressive block expansion \cite{wu2024llama}. To reduce parameter interference, recent techniques employ parameter sparsification, such as TIES-Merging \cite{yadav2023ties} and DARE \cite{yu2024dare}. Other methods address geometric permutation symmetries via optimal transport, coordinate matching, or weight permuting \cite{ainsworth2023rebasin, stoica2023zipit}, perform low-rank decomposition \cite{gargiulo2025svdmerging}, solve closed-form alignments based on activation metadata \cite{jin2023regmean}, perform cross-domain integration \cite{dontuparthi2024cross}, or utilize Fisher information matrices to guide weight aggregation \cite{matena2022fisher, matuskey2024fisher}. Slerp \cite{shoemake1985animating} is also widely used as a geometric baseline for spherical interpolation. However, these methods are static and often suffer from representation misalignment at deeper layers.

\textbf{Representation Calibration.} Representation calibration offers a training-free, feedforward-only alternative to heal representation shift. REPAIR \cite{jordan2022repair} identifies that linear weight interpolation causes "variance collapse"—where intermediate feature standard deviations shrink exponentially layer by layer—and proposes rescaling feature maps to a joint target. However, standard multi-task REPAIR averages statistics across tasks, causing massive cross-task representation interference. To solve this, TCAC \cite{anonymous_tcac} proposes task-conditional calibration using 128 samples per task to record and rescale intermediate activations. While effective, TCAC introduces high procedural complexity with its calibration sets and runtime interception hooks.

\textbf{Test-Time Adaptation (TTA) and BatchNorm.} TTA adjusts a model's parameters or internal states to a target domain using unlabeled test streams \cite{wang2020tent, li2016adabn}. Prior work in TTA has explored source-free hypothesis transfer \cite{liang2020source}, single-sample adaptation \cite{zhang2022memo}, continual test-time adaptation \cite{wang2022continual}, and robust self-training \cite{niu2023towards, niu2022efficient, yuan2023robust}. In model merging, SyMerge \cite{jung2025symerge} and AdaMerging \cite{yang2024adamerging} perform entropy minimization or self-distillation at test time. However, test-time optimization is computationally heavy. Our work is conceptually related to domain-shift adaptation using BatchNorm statistics, such as AdaBN \cite{li2016adabn}, TTN \cite{song2023ttn}, and prediction-time normalization under covariate shift \cite{nado2020evaluating, schneider2020improving}. While those methods adapt standard models to new domains, we apply task-conditional structures to model merging to prevent parameter-level interference from happening in the first place.

\section{Methodology}
\label{method}

Let $\Theta_{\text{pre}}$ denote the weights of a pre-trained base model. We are given $K$ expert models, each independently fine-tuned on a separate task $k \in \{1, \dots, K\}$, resulting in task-specific weights $\Theta_1, \dots, \Theta_K$. Each expert model consists of a shared architecture (such as an encoder or backbone) and a task-specific head (such as a classification layer). We denote the backbone weights as $\theta_k$ and the head weights as $h_k$.

Our goal is to construct a unified multi-task model with a single merged backbone $\theta_{\text{merged}}$ and the original task-specific heads $h_1, \dots, h_K$. During inference on task $k$, the model uses the merged backbone $\theta_{\text{merged}}$ and task head $h_k$ to make predictions:
\begin{equation}
  f_k(x) = h_k(\Phi(x; \theta_{\text{merged}}))
\end{equation}
where $\Phi(x; \theta)$ represents the feature representations output by the backbone.

\subsection{Standard Merging and Variance Collapse}

We consider two standard weight merging baselines:
\begin{enumerate}
    \item \textbf{Weight Averaging (WA):} The merged backbone is computed as the simple element-wise arithmetic mean of the expert backbones:
    \begin{equation}
      \theta_{\text{merged}}^{\text{WA}} = \frac{1}{K} \sum_{k=1}^K \theta_k
    \end{equation}
    \item \textbf{Task Arithmetic (TA):} We extract task vectors $\tau_k = \theta_k - \Theta_{\text{pre}}$ that represent task-specific updates. The merged weights are:
    \begin{equation}
      \theta_{\text{merged}}^{\text{TA}} = \Theta_{\text{pre}} + \lambda \sum_{k=1}^K \tau_k
    \end{equation}
    where $\lambda \in [0, 1]$ is a global scaling coefficient.
\end{enumerate}

When weights are averaged or combined linearly, the resulting activations in deep layers do not scale linearly. Consider a simple linear layer outputting activation $Z = W x$. If we average two independent weight matrices $W_1$ and $W_2$ to form $W_{\text{merged}} = \frac{W_1+W_2}{2}$, the output activation becomes:
\begin{equation}
  Z_{\text{merged}} = W_{\text{merged}} x = \frac{W_1 x + W_2 x}{2}
\end{equation}
If the outputs $W_1 x$ and $W_2 x$ are modeled as independent random variables with variance $\sigma^2$, the variance of the averaged activation $Z_{\text{merged}}$ collapses to:
\begin{equation}
  \text{Var}(Z_{\text{merged}}) = \frac{\text{Var}(W_1 x) + \text{Var}(W_2 x)}{4} = \frac{\sigma^2}{2}
\end{equation}
In deep networks, this variance reduction compounds exponentially layer by layer, resulting in near-zero variance at the final layers. This is known as \textit{variance collapse} \cite{jordan2022repair}, which severely degrades multi-task performance, a phenomenon closely tied to representation dynamics and weight initialization scaling in deep residual networks \cite{he2015delving, he2016identity, glorot2010understanding}.

\begin{theorem}[Variance Collapse and Restoration under TTBN]
\label{thm:variance}
Let $X \in \mathbb{R}^d$ be an input representation. Let $W_1, \dots, W_K \in \mathbb{R}^{d \times d}$ denote the task-specific expert weights, and let $W_{\text{merged}} = \sum_{k=1}^K c_k W_k$ denote the merged weight matrix, where $c_k \ge 0$ and $\sum_{k=1}^K c_k = 1$. Let $Z_k = W_k X$ be the activation under expert $k$, modeled as independent random vectors with zero mean and covariance $\Sigma$. Under weight merging, the uncalibrated merged activation $Z_{\text{merged}} = W_{\text{merged}} X$ has covariance:
\begin{equation}
    \text{Cov}(Z_{\text{merged}}) = \left( \sum_{k=1}^K c_k^2 \right) \Sigma
\end{equation}
For $c_k = 1/K$, this covariance collapses to $\frac{1}{K} \Sigma$. When applied to an intermediate activation vector $Z_{\text{merged}}$ with Channel-wise Test-Time BatchNorm, the resulting normalized activation $\hat{Z}$ exactly restores the target variance:
\begin{equation}
    \text{Var}(\hat{Z}_i) = \gamma_i^2
\end{equation}
where $\gamma_i$ is the learned scaling parameter for channel $i$, completely independent of the weight merging coefficients $c_k$ and the number of tasks $K$.
\end{theorem}

\begin{proof}
By the linearity of expectation, since $\mathbb{E}[Z_k] = 0$, we have:
\begin{equation}
    Z_{\text{merged}} = \sum_{k=1}^K c_k (W_k X) = \sum_{k=1}^K c_k Z_k
\end{equation}
Because the expert activations $Z_1, \dots, Z_K$ are assumed to be independent random vectors, the covariance of their linear combination is:
\begin{align}
    \text{Cov}(Z_{\text{merged}}) &= \mathbb{E}[Z_{\text{merged}} Z_{\text{merged}}^T] \\
    &= \mathbb{E}\left[ \left(\sum_{k=1}^K c_k Z_k\right) \left(\sum_{j=1}^K c_j Z_j\right)^T \right] \\
    &= \sum_{k=1}^K c_k^2 \mathbb{E}[Z_k Z_k^T] = \left(\sum_{k=1}^K c_k^2 \right) \Sigma
\end{align}
If $c_k = 1/K$, then $\sum_{k=1}^K c_k^2 = \sum_{k=1}^K \frac{1}{K^2} = \frac{1}{K}$. Thus, $\text{Cov}(Z_{\text{merged}}) = \frac{1}{K}\Sigma$. For a network with $L$ independent sequential layers, the variance shrinks exponentially as $K^{-L}$, resulting in catastrophic variance collapse as $L \to \infty$.

Now, let $Z_{b, i}$ denote the activation of channel $i$ for sample $b$ in a test batch of size $B$. Test-Time BatchNorm computes the exact sample mean $\mu_i = \frac{1}{B}\sum_b Z_{b, i}$ and sample variance $\sigma_i^2 = \frac{1}{B}\sum_b (Z_{b, i} - \mu_i)^2$ on the fly. The normalized output is:
\begin{equation}
    \hat{Z}_{b, i} = \gamma_i \left( \frac{Z_{b, i} - \mu_i}{\sqrt{\sigma_i^2 + \epsilon}} \right) + \beta_i
\end{equation}
Taking the sample variance of the normalized activation $\hat{Z}_i$ across the batch yields:
\begin{align}
    \text{Var}(\hat{Z}_i) &= \frac{1}{B} \sum_{b=1}^B (\hat{Z}_{b, i} - \bar{\hat{Z}}_i)^2 \\
    &= \frac{1}{B} \sum_{b=1}^B \left( \gamma_i \left( \frac{Z_{b, i} - \mu_i}{\sqrt{\sigma_i^2 + \epsilon}} \right) + \beta_i - \beta_i \right)^2 \\
    &= \gamma_i^2 \left( \frac{\frac{1}{B}\sum_{b=1}^B (Z_{b, i} - \mu_i)^2}{\sigma_i^2 + \epsilon} \right) \approx \gamma_i^2
\end{align}
where the approximation becomes an exact equality in the limit as $\epsilon \to 0$. This completes the proof.
\end{proof}

\subsection{Proposed Method: Test-Time BatchNorm (TTBN)}
\label{sec:ttbn}

To resolve variance collapse without the need for calibration datasets, forward passes, or runtime interception hooks, we propose \textbf{Test-Time BatchNorm (TTBN)}.

TTBN is built on a simple, minimalist realization: the collapsed running statistics stored in the merged BatchNorm layers are what drive representation distortion. Instead of using those static, incorrect running stats, we should use the actual, dynamic statistics computed from the test data itself on-the-fly.

Specifically, TTBN performs standard weight merging on all parameters of the backbone. During evaluation, for any intermediate layer with BatchNorm \cite{ioffe2015batch}, we compute the batch mean and batch variance directly from the incoming activation tensor $X \in \mathbb{R}^{B \times C \times H \times W}$ across the spatial and batch dimensions:
\begin{align}
    \mu_{\text{batch}, c} &= \frac{1}{B \cdot H \cdot W} \sum_{b=1}^B \sum_{h=1}^H \sum_{w=1}^W X_{b, c, h, w} \\
    \sigma^2_{\text{batch}, c} &= \frac{1}{B \cdot H \cdot W} \sum_{b=1}^B \sum_{h=1}^H \sum_{w=1}^W (X_{b, c, h, w} - \mu_{\text{batch}, c})^2
\end{align}
The normalized activation is then computed on-the-fly as:
\begin{equation}
    \hat{X}_{b, c, h, w} = \frac{X_{b, c, h, w} - \mu_{\text{batch}, c}}{\sqrt{\sigma^2_{\text{batch}, c} + \epsilon}} \cdot \gamma_c + \beta_c
\end{equation}
where $\gamma_c$ and $\beta_c$ are the learned scale and bias parameters of the BatchNorm layer.

In PyTorch, this is implemented extremely elegantly by setting the model to training mode during evaluation (`model.train()`) but freezing all running statistic updates by setting `momentum = 0.0`. This ensures that:
\begin{enumerate}
    \item Activations are normalized using the exact mean and variance of the current test batch.
    \item The model's stored running statistics are not corrupted by the test data.
\end{enumerate}

While alternative normalization techniques like Layer Normalization \cite{ba2016layer}, Instance Normalization \cite{ulyanov2016instance}, Group Normalization \cite{wu2018group}, Weight Normalization \cite{salimans2016weight}, and Spectral Normalization \cite{miyato2018spectral} do not rely on running statistics (and thus are naturally immune to static variance collapse), convolutional backbones like ResNet-18 are heavily built on BatchNorm, and its optimization landscape is deeply dependent on it \cite{santurkar2018how}. Thus, TTBN offers a highly effective way to salvage standard BatchNorm networks without modifying their architecture, and requires zero extra parameters, zero task-conditional routing, and zero offline setup, adapting dynamically and perfectly to any feature scale or representation shift in a zero-shot manner.

\subsection{Sequential Test-Time BatchNorm (STTBN)}
\label{sec:sttbn}

A potential vulnerability of standard TTBN is its dependency on test-time batch size. When the batch size $B$ is very small (e.g., $B = 2$ or $B = 4$), the sample-based statistics $\mu_{\text{batch}, c}$ and $\sigma^2_{\text{batch}, c}$ exhibit extremely high variance, leading to noisy normalization that can degrade downstream accuracy. 

To overcome this limitation in a purely online, minimalist manner, we propose \textbf{Sequential Test-Time BatchNorm (STTBN)}. Instead of treating each test batch as an isolated sample, STTBN maintains and sequentially updates running estimates of the test-time mean ($\hat{\mu}_{\text{running}}$) and variance ($\hat{\sigma}^2_{\text{running}}$) as the test stream progresses, utilizing a momentum parameter $\beta \in (0, 1]$:
\begin{align}
    \hat{\mu}_{\text{running}, c}^{(t)} &= (1 - \beta) \hat{\mu}_{\text{running}, c}^{(t-1)} + \beta \mu_{\text{batch}, c}^{(t)} \\
    \hat{\sigma}^{2(t)}_{\text{running}, c} &= (1 - \beta) \hat{\sigma}^{2(t-1)}_{\text{running}, c} + \beta \sigma^{2(t)}_{\text{batch}, c}
\end{align}
To normalize the activations of the current batch on-the-fly, STTBN can operate in one of two configurations:
\begin{enumerate}
    \item \textbf{Pre-Hook Normalization:} The running statistics are updated \textit{before} the current layer's forward pass, integrating the current batch's scale to stabilize normalization instantly:
    \begin{equation}
        \hat{X}_{b, c, h, w}^{(t)} = \frac{X_{b, c, h, w}^{(t)} - \hat{\mu}_{\text{running}, c}^{(t)}}{\sqrt{\hat{\sigma}^{2(t)}_{\text{running}, c} + \epsilon}} \cdot \gamma_c + \beta_c
    \end{equation}
    \item \textbf{Post-Hook Normalization:} The current batch is normalized using historical statistics $\hat{\mu}_{\text{running}, c}^{(t-1)}$ and $\hat{\sigma}^{2(t-1)}_{\text{running}, c}$, after which the statistics are updated for the next batch.
\end{enumerate}

By accumulating statistics sequentially, STTBN filters out single-batch statistical noise, enabling highly robust, stable, and accurate normalization even for batch sizes as small as 2. We initialize the running statistics using either the merged model's static running statistics (Merged Init) or the task-specific expert's native statistics (Expert Init), with the latter providing a well-scaled starting point that further improves sample-efficient adaptation.
\subsection{Task-Specific BatchNorm Merging (TSBM)}

As a structural decoupling baseline, we also evaluate \textbf{Task-Specific BatchNorm Merging (TSBM)}. In TSBM, we partition the backbone weights $\mathcal{W}$ into two disjoint sets:
\begin{equation}
    \mathcal{W} = \mathcal{W}_{\text{conv/linear}} \cup \mathcal{W}_{\text{BN}}
\end{equation}
where $\mathcal{W}_{\text{conv/linear}}$ contains all parameters belonging to convolutional and linear layers, and $\mathcal{W}_{\text{BN}}$ contains all parameters belonging to BatchNorm layers (including running statistics and affine parameters).

In TSBM, we merge only the convolutional and linear parameters using standard weight merging (such as Weight Averaging or Task Arithmetic):
\begin{equation}
    \mathcal{W}_{\text{conv/linear}}^{\text{merged}} = \text{Merge}(\{\mathcal{W}_{\text{conv/linear}, k}\}_{k=1}^K)
\end{equation}
However, we do not merge the BatchNorm parameters. Instead, we keep them task-conditional. For each task $k$, we construct a task-specific merged model backbone:
\begin{equation}
    \theta_{\text{TSBM}}^{(k)} = \mathcal{W}_{\text{conv/linear}}^{\text{merged}} \cup \mathcal{W}_{\text{BN}, k}
\end{equation}
where $\mathcal{W}_{\text{BN}, k}$ represents the native BatchNorm parameters and statistics of Expert $k$. During inference on task $k$, we swap in the unmerged BatchNorm parameters and statistics of Expert $k$. 

While TSBM completely eliminates the need for calibration datasets and runtime hooks, we show in Section~\ref{analysis} that it suffers from a critical limitation: a severe scale mismatch between the static unmerged running statistics and the collapsed feature scales of the merged backbone.

\subsection{Demystifying and Correcting Prior Hook-Based Calibration}

Prior frameworks like TCAC and REPAIR attempt to correct representation shift by recording intermediate activation mean and variance statistics on a calibration set of $N=128$ samples, and then registering forward hooks to rescale activations at inference time:
\begin{equation}
    Y_{\text{calibrated}} = \frac{Y - \mu_{\text{merged}}}{\sigma_{\text{merged}}} \cdot \sigma_{\text{orig}} + \mu_{\text{orig}}
\end{equation}

However, we identify a fundamental methodological flaw in how these statistics are collected. In previous implementations, the statistics $\mu_{\text{merged}}$ and $\sigma_{\text{merged}}$ are collected *simultaneously* across all layers of the uncalibrated merged model. 

In a deep neural network, the activation $X^{(l+1)}$ entering layer $l+1$ is a function of the output $Y^{(l)}$ of layer $l$. If we do not calibrate layer $l$, the statistics of $X^{(l+1)}$ are highly distorted. When we collect $\mu_{\text{merged}}^{(l+1)}$ and $\sigma_{\text{merged}}^{(l+1)}$ on the uncalibrated model, they represent this distorted distribution. 

However, during evaluation, we register calibration hooks on *all* layers. This means that when the forward pass executes, layer $l$ *is* being calibrated. Thus, the activation entering layer $l+1$ has been restored to a sane distribution. But the hook at layer $l+1$ still subtracts $\mu_{\text{merged}}^{(l+1)}$ and divides by $\sigma_{\text{merged}}^{(l+1)}$—statistics that were collected on the *uncalibrated* distorted inputs! 

Applying this mismatched calibration transformation catastrophically destroys the signals in the network, leading to random-guessing performance ($\sim$11\% accuracy) as shown in Section~\ref{exper}.

To fix this, we implement a mathematically sound, sequential Offline BN Recalibration baseline (Corrected REPAIR/TCAC). Instead of registering concurrent hooks, we set the merged model to training mode (`model.train()`) with frozen weights and run forward passes on the 128 calibration samples with `momentum = 1.0`. This sequentially and exactly re-estimates the running mean and variance of every BatchNorm layer in a feedforward manner, taking into account the calibrated outputs of all preceding layers. 

\section{Experiments}
\label{exper}

We evaluate our proposed methods on three standard image classification datasets: MNIST \cite{lecun1998mnist}, Fashion-MNIST \cite{xiao2017fashion}, and CIFAR-10 \cite{krizhevsky2009cifar}. All grayscale images (MNIST, Fashion-MNIST) are converted to 3-channel RGB and resized to $32 \times 32$ pixels.

\subsection{Setup and Hyperparameters}

We employ a standard **ResNet-18** backbone pre-trained on ImageNet. Each expert model is fine-tuned independently on its respective task for 5 epochs using AdamW with a learning rate of $5 \times 10^{-4}$ and weight decay of $10^{-4}$, with a batch size of 128.

For REPAIR (Simultaneous) and TCAC (Simultaneous), we use a random calibration set of $N=128$ samples from each training set. For TSBM, we merge Conv and Linear weights using the best-performing lambda from Task Arithmetic. For our sequential Offline BN Recalibration (REPAIR-Corrected) baseline, we run a single forward pass with `momentum = 1.0` on the same 128-sample calibration set.

\subsection{Main Results}

Table~\ref{tab:results} summarizes the classification accuracies across all methods.

\begin{table*}[t]
\caption{Multi-task model merging accuracies (\%) on ResNet-18. Best merging results are highlighted in \textbf{bold}.}
\label{tab:results}
\vskip 0.15in
\begin{center}
\begin{small}
\begin{tabular}{lcccc}
\toprule
Method & MNIST & Fashion-MNIST & CIFAR-10 & Average \\
\midrule
Individual Experts (Upper Bound) & 99.30\% & 91.38\% & 81.31\% & 90.66\% \\
\midrule
Weight Averaging (WA) & 21.82\% & 10.27\% & 21.96\% & 18.02\% \\
WA + TTBN (Ours) & 96.83\% & 85.38\% & 62.05\% & 81.42\% \\
Task Arithmetic (TA) ($\lambda = 0.2$) & 42.64\% & 43.32\% & 30.25\% & 38.74\% \\
\textbf{TA + TTBN (Ours, Proposed)} ($\lambda = 0.4$) & \textbf{97.31\%} & \textbf{86.02\%} & \textbf{62.52\%} & \textbf{81.95\%} \\
\midrule
REPAIR (Simultaneous, Broken) & 11.35\% & 10.00\% & 12.68\% & 11.34\% \\
TCAC (Simultaneous, Broken) & 11.35\% & 10.00\% & 12.68\% & 11.34\% \\
REPAIR (Offline Recalibration, Corrected Baseline) & 93.62\% & 81.21\% & 53.73\% & 76.19\% \\
\midrule
TSBM (Ours, Decoupled Baseline) ($\lambda = 0.7$) & 24.42\% & 37.45\% & 30.82\% & 30.90\% \\
\bottomrule
\end{tabular}
\end{small}
\end{center}
\vskip -0.1in
\end{table*}

\section{Analysis and Discussion}
\label{analysis}

\subsection{Ablation of Merging Coefficient ($\lambda$)}

Figure~\ref{fig:sweep} shows the performance of Task Arithmetic, TA + TTBN, and TSBM across different values of the merging coefficient $\lambda$.

\begin{figure}[h]
\vskip 0.2in
\begin{center}
\includegraphics[width=\columnwidth]{merging_performance_sweep.png}
\caption{Performance comparison of model merging methods across different values of the merging coefficient $\lambda$.}
\label{fig:sweep}
\end{center}
\vskip -0.2in
\end{figure}

\subsection{Robustness under Small Batch Sizes: Sequential Test-Time BatchNorm (STTBN)}
\label{sec:exp_sttbn}

An important practical consideration of Test-Time BatchNorm (TTBN) is its reliance on the batch size of the test stream. Because TTBN computes mean and variance statistics on-the-fly during inference, extremely small batch sizes may lead to noisy statistical estimates, which can degrade classification accuracy.

To analyze and completely resolve this sensitivity in a purely online, training-free manner, we evaluate our proposed \textbf{Sequential Test-Time BatchNorm (STTBN)} across a range of test-time batch sizes from $2$ to $128$ under different momentum values $\beta \in \{0.01, 0.05, 0.1, 0.2\}$. Table~\ref{tab:sttbn_robustness} summarizes these results.

\begin{table*}[t]
\caption{Multi-task model merging average classification accuracies (\%) on ResNet-18 across different test-time batch sizes and momentum values ($\beta$) under our Sequential Test-Time BatchNorm (STTBN) framework. Best results per batch size are highlighted in \textbf{bold}.}
\label{tab:sttbn_robustness}
\vskip 0.15in
\begin{center}
\begin{small}
\setlength{\tabcolsep}{4.5pt}
\begin{tabular}{lccccccc}
\toprule
Method / Momentum ($\beta$) & BS = 2 & BS = 4 & BS = 8 & BS = 16 & BS = 32 & BS = 64 & BS = 128 \\
\midrule
Baseline TTBN (No Sequential Updates) & 43.75\% & 61.45\% & 69.91\% & 76.39\% & 79.84\% & 81.16\% & \textbf{81.95}\% \\
\midrule
STTBN (Pre-Hook, Merged Init, $\beta = 0.01$) & \textbf{80.99}\% & 80.35\% & 78.67\% & 75.15\% & 68.52\% & 54.93\% & 35.30\% \\
STTBN (Pre-Hook, Merged Init, $\beta = 0.05$) & 78.69\% & \textbf{80.58}\% & \textbf{80.90}\% & 80.56\% & 79.35\% & 76.73\% & 71.57\% \\
STTBN (Pre-Hook, Merged Init, $\beta = 0.1$) & 74.44\% & 78.95\% & 80.50\% & \textbf{80.90}\% & 80.68\% & 79.52\% & 76.99\% \\
STTBN (Pre-Hook, Merged Init, $\beta = 0.2$) & 66.63\% & 75.33\% & 79.07\% & 80.61\% & \textbf{81.25}\% & \textbf{80.84}\% & 79.90\% \\
\bottomrule
\end{tabular}
\end{small}
\end{center}
\vskip -0.1in
\end{table*}

We observe three crucial scientific insights from Table~\ref{tab:sttbn_robustness}:
\begin{enumerate}
    \item \textbf{Catastrophic Noise Mitigated:} Standard TTBN performs poorly at batch size 2 ($43.75\%$) due to extreme single-batch statistical noise. By introducing sequential updates, STTBN with $\beta = 0.01$ achieves an outstanding \textbf{80.99\%} average accuracy, completely bridging the gap to the large-batch upper bound ($81.95\%$) and outperforming standard TTBN by a massive \textbf{+37.24\%} absolute!
    \item \textbf{Batch Size vs. Momentum Trade-off:} There is an elegant, inverse relationship between the batch size and the optimal momentum parameter $\beta$. When the batch size is extremely small (e.g., $BS=2$), the test stream consists of a large number of batches (5,000 steps). Therefore, a very small momentum ($\beta = 0.01$) is optimal as it slowly and smoothly filters out single-batch noise. Conversely, when the batch size is large (e.g., $BS=32$ or $64$), the stream has fewer steps (only 312 or 156 steps), requiring a larger momentum ($\beta = 0.2$) or standard TTBN ($\beta = 1.0$) to adapt rapidly.
    \item \textbf{Pre-Hook Superiority:} We also compared Pre-Hook updates (where stats are updated before normalization) with Post-Hook updates (where stats are updated after normalization using historical statistics). Pre-Hook sequential updates are strictly superior to Post-Hook updates (e.g., 66.63\% vs 60.09\% for BS=2, $\beta=0.2$) because incorporating the current batch's statistics on-the-fly immediately scales and stabilizes representations.
\end{enumerate}

These results prove that our sequential dynamic normalization framework is incredibly robust to batch-size constraints. It completely eliminates the need for offline calibration sets or task-specific metadata storage, offering an optimal, purely online adaptation solution for real-world model merging.

\subsection{Why TSBM Underperforms}

While Task-Specific BatchNorm Merging (TSBM) keeps the expert's unmerged BatchNorm parameters and running statistics, it fails to completely recover performance (achieving only 30.90\%). This is because the convolutional weights in the backbone are merged, which inevitably shrinks the magnitude of their features.

When these shrunken representations are passed through unmerged expert BN layers, we subtract $\mu_{\text{orig}}$ and divide by the original, large uncollapsed running standard deviation $\sigma_{\text{orig}}$. This division by a mismatched, large variance scales the already-shrunken activations down even further, exponentially dampening features and leading to representation collapse in deep layers. This demonstrates that static unmerged running statistics are fundamentally incompatible with merged backbone weights.

\subsection{Why TTBN and Corrected Offline Calibration Excel}

In contrast, both TTBN (Ours) and Offline BN Recalibration (Corrected REPAIR) resolve this mismatch by dynamically or sequentially updating the running statistics to match the actual merged representations. 

By running forward passes on a calibration set, Offline BN Recalibration sequentially propagates activations, ensuring that the running mean and variance computed at layer $l+1$ are based on the correctly calibrated outputs of layer $l$. This sequential feedforward updates perfectly capture the shrunken scales, achieving a high average accuracy of 76.19\%.

Our proposed Test-Time BatchNorm (TTBN) takes this a step further. Instead of relying on an offline calibration dataset of 128 samples (which introduces data dependence and requires a separate feedforward run), TTBN computes these exact activation scales directly from the incoming test stream on the fly. This zero-shot adaptation achieves 81.95\% average accuracy, outperforming even the corrected offline baseline on the harder CIFAR-10 dataset (62.52\% vs 53.73\%). This proves that dynamic, test-time normalization is the most optimal, elegant, and minimalist solution for multi-task model merging.

\section{Conclusion}
\label{conclusion}

In this paper, we challenge the growing complexity of training-free activation calibration methods in model merging. Guided by Occam's razor, we show that simply computing batch statistics on-the-fly via Test-Time BatchNorm (TTBN) completely resolves activation variance collapse and representation shift. TTBN achieves state-of-the-art results on multi-task model merging without any calibration data, calibration forward passes, or runtime interception hooks. We recommend that the community shift focus towards simple, dynamic test-time adaptations over mathematically convoluted offline correction pipelines.

\bibliography{submission}
\bibliographystyle{icml2026}

\end{document}
